\documentclass{article}
\usepackage{amsmath,amssymb}
\usepackage{xurl}
\usepackage[hidelinks]{hyperref}
% Shared commands for notes in docs/paper-gaps/.

\DeclareUnicodeCharacter{2080}{\ifmmode{}_{0}\else\textsubscript{0}\fi}
\DeclareUnicodeCharacter{2081}{\ifmmode{}_{1}\else\textsubscript{1}\fi}
\DeclareUnicodeCharacter{2082}{\ifmmode{}_{2}\else\textsubscript{2}\fi}
\DeclareUnicodeCharacter{2099}{\ifmmode{}_{n}\else\textsubscript{n}\fi}

\newcommand{\Lean}{\textsc{Lean}}
\ExplSyntaxOn
\NewDocumentCommand{\leanid}{m}{
  \ifmmode
    \textsf{\detokenize{#1}}
  \else
    \group_begin:
    \urlstyle{sf}
    \tl_set:Nn \l_tmpa_tl {#1}
    \tl_replace_all:Nnn \l_tmpa_tl {\_} {_}
    \exp_args:NV \path \l_tmpa_tl
    \group_end:
  \fi
}
\ExplSyntaxOff
\providecommand{\tr}{\operatorname{tr}}
\newcommand{\tnleanIssueBase}{https://github.com/LionSR/TNLean/issues/}
\newcommand{\tnleanPrBase}{https://github.com/LionSR/TNLean/pull/}
\newcommand{\tnleanDocBase}{https://sirui-lu.com/TNLean/docs/}
\newcommand{\ghissue}[1]{\href{\tnleanIssueBase#1}{GitHub issue \##1}}
\newcommand{\ghpr}[1]{\href{\tnleanPrBase#1}{PR \##1}}
\newcommand{\doclink}[1]{\href{\tnleanDocBase#1.html}{\path{#1}}}

% Dirac notation and the identity operator, as in blueprint/src/macros/common.tex.
% \providecommand keeps note-local definitions authoritative.
\usepackage{dsfont}
\providecommand{\ket}[1]{|#1\rangle}
\providecommand{\bra}[1]{\langle #1|}
\providecommand{\braket}[2]{\langle #1 | #2 \rangle}
\providecommand{\ketbra}[2]{|#1\rangle\!\langle #2|}
\providecommand{\Id}{\mathds{1}}
\providecommand{\one}{\mathds{1}}
\providecommand{\C}{\mathbb{C}}
\providecommand{\MN}[1]{M_{#1}(\C)}
\providecommand{\MD}{M_D(\C)}

% Verdict marker: \gapnote{<kind>}{<status>}. Kind is the policy
% classification (clarification, local-correction, scope-restriction,
% unfaithful, false-source, open-gap); status is open, wip (elimination
% underway), resolved, or historical. Machine-read by the site index;
% prints a verdict line.
\newcommand{\gapnote}[2]{\par\noindent\textbf{Verdict:} #1 (#2).\par}

\title{Exact Preparation at a Zero Subleading Transfer Spectrum}
\author{}
\date{2026-10-03}
\begin{document}
\maketitle
\gapnote{local-correction}{resolved}

\section{The source prescription}

For a normal tensor in trace-preserving gauge, the source writes
$\mathcal E=P+R$, where $P(X)=\tr(X)\sigma$, $\sigma>0$,
$\tr\sigma=1$, and $R$ has spectral radius less than one
\cite[paragraph ``Preliminaries'']{Malz2024Preparation}.
Its logarithmic block prescription uses the correlation length
$\xi=-1/\log|\lambda_2|$ and chooses
$q=\lceil2\xi(1+\eta)\log N\rceil$ after Lemma 1.
When all subleading eigenvalues vanish, $\log|\lambda_2|$ is not finite.
The prescription must therefore be replaced by a positive finite block
length at this endpoint. In particular, a value assigned to the symbol
$\xi$ alone is not the spectral hypothesis used below.

\section{Nilpotent transients and a finite block length}

Assume explicitly that every eigenvalue of $\mathcal E$ distinct from
$1$ is zero. Normality makes the leading fixed projection simple, and
trace preservation and stationarity give $P^2=P$ and $PR=RP=0$.
Every eigenvalue of $R$ is zero. On the $D^2$-dimensional matrix space,
the generalized eigenspace decomposition gives $R^{D^2}=0$, including
nondiagonalizable maps. Consequently, for every $n\geq D^2$,
\begin{align}
    \mathcal E^n(X)&=\tr(X)\sigma.
    \label{eq:mswc24_zero_spectrum_transfer}
\end{align}
The subleading part need not vanish before that finite blocking length.

For a block $B$ of at least $D^2$ sites, the transfer/polar relation of
the source gives $\widehat B^\dagger\widehat B=\sigma^{\mathsf T}\otimes I$.
Thus its positive polar factor is exactly
$(\sqrt\sigma)^{\mathsf T}\otimes I$, and it is injective because
$\sigma>0$. Replacing the polar factors by the fixed-point factors in
the source construction therefore changes no state.

Choose $q_0=D^2+3D+1$. For every $N\geq q_0$, write
$N=Mq_0+r$ with $M\geq1$ and $0\leq r<q_0$.
Use $M-1$ blocks of length $q_0$ and one of length $q_0+r$.
All lengths lie between $q_0$ and $2q_0$, so the local-isometry
preparation applies in depth at most $2C(d,D)q_0$.
The raw periodic vector already has norm one, and the circuit prepares
its normalized vector exactly. This bound is independent of $N$ and
requires no divisibility condition.

\section{Formal scope}

The construction is the exact-factor endpoint of the source equations
\texttt{eq:B\_TM}, \texttt{eq:phi\_tilde}, and
\texttt{eq:app\_tidle\_phi\_1}; it is not an evaluation of the printed
logarithmic prescription at zero correlation length. Its spectral
assumption allows nilpotent Jordan blocks and supplies no separate
injectivity or diagonalizability certificate.
The preparation theorem is
\leanid{MPSPreparation.exists_isPreparedInDepth_normalizedMPVState_of_subleading_eigenvalues_eq_zero}.
It fixes $d,D$ before choosing the depth and applies to every normal
left-canonical tensor and normalized positive definite fixed state
satisfying the stated zero-subleading-spectrum condition.

This endpoint consequence does not settle the polynomial-accuracy
construction at positive correlation length or its arbitrary-length
extensions. Those results have their own hypotheses and proofs.

\bibliographystyle{alpha}
\bibliography{references}
\end{document}
